\documentclass[12pt,a4paper]{article}
\usepackage[margin=25mm, headheight=15pt, headsep=10pt]{geometry}
\usepackage{fontspec}
\usepackage[table]{xcolor} % Note: table option is required for rowcolors
\usepackage{titlesec}
\usepackage{tocloft}
\usepackage{fancyhdr}
\usepackage{booktabs}
\usepackage{tcolorbox}
\tcbuselibrary{skins, breakable}
\usepackage[colorlinks=true, linkcolor=emerald, urlcolor=emerald, bookmarks=true]{hyperref}

\definecolor{emerald}{RGB}{0,102,51}
\definecolor{darkred}{RGB}{139,0,0}
\definecolor{lightgray}{RGB}{245,245,245}

\usepackage[nil]{babel}
\babelprovide[import, main]{bengali}
\babelprovide[import]{arabic}
\babelfont{rm}[Renderer=HarfBuzz,Scale=1.0]{Noto Serif Bengali}
\babelfont{sf}[Renderer=HarfBuzz]{Noto Sans Bengali}
\babelfont{tt}[Renderer=HarfBuzz]{Noto Sans Bengali}
\babelfont[arabic]{rm}[Renderer=HarfBuzz,Scale=1.1]{Noto Naskh Arabic}

% Section styling
\titleformat{\section}{\Large\bfseries\color{emerald}}{\thesection.}{0.5em}{}
\titleformat{\subsection}{\large\bfseries\color{emerald!85!black}}{\thesubsection.}{0.5em}{}
\titleformat{\subsubsection}{\normalsize\bfseries\color{emerald!70!black}}{\thesubsubsection.}{0.5em}{}
\titlespacing*{\section}{0pt}{18pt}{8pt}

% Fancy Header/Footer setup
\pagestyle{fancy}
\fancyhf{} % clear all header and footer fields
\fancyhead[L]{\color{emerald}\small\textbf{\nouppercase{\leftmark}}}
\fancyfoot[L]{\scriptsize\color{black!50}{\lat{AI}}-সংকলিত, আলেম-যাচাইকৃত নয়}
\fancyfoot[C]{\color{emerald!70!black}\thepage}
\renewcommand{\headrulewidth}{0.4pt}
\renewcommand{\headrule}{\hbox to\headwidth{\color{emerald}\leaders\hrule height \headrulewidth\hfill}}
\renewcommand{\footrulewidth}{0pt}

% Custom boxes
% 1. Ayat/Hadith Box
\newtcolorbox{ayatbox}[1][]{
    enhanced, breakable, 
    colback=emerald!5, 
    colframe=emerald, 
    boxrule=0.5pt, 
    arc=3pt, 
    left=10pt, right=10pt, top=10pt, bottom=10pt,
    #1
}

% 2. Quote/Translation Box
\newtcolorbox{quotebox}[1][]{
    enhanced, breakable,
    frame hidden,
    colback=lightgray,
    borderline west={3pt}{0pt}{emerald},
    left=15pt, right=10pt, top=10pt, bottom=10pt,
    sharp corners,
    #1
}

% 3. AI-generated notice box
\newtcolorbox{warnbox}[1][]{
    enhanced, breakable,
    colback=red!4,
    colframe=darkred,
    boxrule=0.8pt,
    arc=3pt,
    left=10pt, right=10pt, top=10pt, bottom=10pt,
    #1
}

% Commands
\newcommand{\arab}[1]{\foreignlanguage{arabic}{#1}}
\newcommand{\kib}[1]{\textcolor{emerald}{\textbf{#1}}}

% Latin (English) fallback: Noto Serif Bengali has NO Latin glyphs
\newfontfamily\latfont[Renderer=HarfBuzz]{Noto Serif}
\newcommand{\lat}[1]{{\latfont #1}}

\setlength{\parskip}{5pt}
\setlength{\parindent}{0pt}

% Fix for missing bullet in Noto Serif Bengali
\renewcommand{\labelitemi}{$\bullet$}



\begin{document}
\begin{center}{\Large\bfseries\color{emerald} ঘটনা ১ — নবীজি (সা.): তাহাজ্জুদে পা ফুলে যাওয়া — 'কৃতজ্ঞ বান্দা'}\end{center}
\vspace{1em}
\section*{ঘটনা ১ — নবীজি (সাল্লাল্লাহু আলাইহি ওয়াসাল্লাম): তাহাজ্জুদে পা ফুলে যাওয়া — 'কৃতজ্ঞ বান্দা'}

\begin{warnbox}
 \textbf{গুরুত্বপূর্ণ নোটিশ (\lat{Important} \lat{Notice}):} এই কন্টেন্টটি \lat{AI} (কৃত্রিম বুদ্ধিমত্তা) দিয়ে প্রস্তুত ও সংকলিত। \lat{AI} কঠোর সূত্র-মান নীতি মেনে শুধু নির্ভরযোগ্য উৎস থেকে তথ্য সংগ্রহ করেছে — কেবল সহীহ/হাসান হাদিস ও সহীহ সনদের আসার রাখা হয়েছে, দুর্বল সনদ বাদ দেওয়া হয়েছে। তবে এটি কোনো আলেম বা মুহাদ্দিস কর্তৃক যাচাইকৃত নয়।
\end{warnbox}

\begin{quote}
\textbf{সম্পর্কিত ফাইল:} পার্ট ১ (মর্যাদা), পার্ট ৫ঘ (তাহাজ্জুদ-নফল)।
\end{quote}

\medskip\hrule\medskip

\subsection*{১. সূত্র কার্ড (\lat{Source} \lat{Card})}

\begin{table}[h]\centering\small
\begin{tabular}{ll}
বিষয় & বিবরণ \\
\hline
\textbf{ঘটনা} & নবীজি (সা.)-এর দীর্ঘ তাহাজ্জুদে পা ফুলে যাওয়া \\
\textbf{রাবি} & আয়েশা (রাঃ) \\
\textbf{গ্রন্থ ও নম্বর} & সহীহ বুখারি ১১৩০; সহীহ মুসলিম ৭৪১ \\
\textbf{সনদের মান} & সহীহ (সহীহাইন) — সর্বোচ্চ মানের দলিল \\
\hline
\end{tabular}
\end{table}

\medskip\hrule\medskip

\subsection*{২. পটভূমি (\lat{Background})}

\textbf{কে:} নবীজি মুহাম্মাদ (সা.); আয়েশা (রাঃ) — তাঁর স্ত্রী ও ঘনিষ্ঠ সহচরী।
\textbf{কখন:} মদিনার মদিনা যুগ — রাতের ইবাদতের নিয়মিত অভ্যাস। বিশেষ কোনো সময় বাঁধা ছিল না; বছরের যেকোনো রাতে।
\textbf{কোথায়:} নবীজির ঘর — আয়েশা (রাঃ)-এর কাছে।
\textbf{কেন:} কৃতজ্ঞতার প্রশিক্ষণ ও আল্লাহর কাছে প্রিয় বান্দা হওয়ার গভীর পিপাসা — "তাহাজ্জুদ" নফল নামাজের মাধ্যমে।

\medskip\hrule\medskip

\subsection*{৩. পূর্ণ ঘটনার বর্ণনা (\lat{The} \lat{Full} \lat{Event})}

আয়েশা (রাঃ) বলেন: নবীজি (সা.) রাতের নামাজ (তাহাজ্জুদ) পড়তেন — ততক্ষণ পর্যন্ত দাঁড়িয়ে, \textbf{যতক্ষণ তাঁর পা ফুলে চিড় না হয়ে যেত}।

একবার আয়েশা (রাঃ) সবিস্ময়ে বললেন:
\begin{quote}
"হে আল্লাহর রাসূল! আপনি এত কষ্ট করেন অথচ আল্লাহ আপনার আগের-পরের সমস্ত গুনাহ মাফ করে দিয়েছেন — কেন তবু এত (দীর্ঘ) ইবাদত?"
\end{quote}

নবীজি (সা.) সুললিত কণ্ঠে জবাব দিলেন:
\begin{quote}
\textbf{\lat{“}আফালা আকূনু 'আবদান শাকুরা — আমি কি কৃতজ্ঞ বান্দা হব না?\lat{”}}
\end{quote}

অর্থাৎ — "আমি কি এমন বান্দা হব না যে সেই রবের কৃতজ্ঞতা আদায় করে, যিনি আমার সব গুনাহ মাফ করে দিয়েছেন? মাফ পাওয়ার কারণে কি ইবাদত কমে যায়? বরং বেড়ে যায় — কেননা এটাই কৃতজ্ঞতার প্রকাশ।"

\textbf{শব্দগত বিশ্লেষণ:} \lat{“}শাকুরা\lat{”} শব্দটি \lat{“}শাকির\lat{”}-থেকে বেশি গভীর — শাকুর সেই কৃতজ্ঞ, যে সামান্য অনুগ্রহেও অধিক, নিরবচ্ছিন্ন ও স্থায়ী কৃতজ্ঞতা প্রকাশ করে — শুধু মুখে নয়, অঙ্গ-প্রত্যঙ্গের আমলে।

\medskip\hrule\medskip

\subsection*{৪. ধাপে ধাপে বিশ্লেষণ (\lat{Step-by-Step} \lat{Analysis})}

\begin{enumerate}
\item \textbf{আয়েশা (রাঃ)-এর প্রশ্নের পেছনের যুক্তি:} সাধারণ মানুষ মনে করে — দোষ-ত্রুটি মাফ করলে ইবাদতের দায়িত্ব কমে যায়। আয়েশা (রাঃ) সেই প্রাকৃতিক প্রশ্নই করলেন — মানুষের দুর্বল দৃষ্টিভঙ্গি প্রকাশ পায়।
\item \textbf{নবীজির (সা.) উত্তর:} ইবাদতের উদ্দেশ্য \textbf{শাস্তির ভয় নয়}, \textbf{কৃতজ্ঞতার আদায়}। যত বেশি অনুগ্রহ, তত বেশি কৃতজ্ঞ বান্দা হওয়া উচিত।
\item \textbf{পা ফুলে যাওয়ার বাস্তবতা:} রাতের নামাজ কত দীর্ঘ হতো — দাড়িয়ে দাড়িয়ে পা ফুলে যেত। এটি \textbf{শারীরিক ক্লান্তির বিনিময়ে আধ্যাত্মিক তৃপ্তির} মূর্ত প্রতীক।
\item \textbf{মোকাবিলার শিক্ষা:} নবীজি (সা.)-এর শারীরিক অবস্থা তাঁকে ইবাদত থেকে বিরত করেনি — এটি সাহাবাদের জন্য রোল-মডেল।
\end{enumerate}
\medskip\hrule\medskip

\subsection*{৫. শব্দের ব্যাখ্যা (\lat{Key} \lat{Words})}

\begin{itemize}
\item \textbf{কৃতজ্ঞতা (\lat{Shukr})}: অনুগ্রহকে স্বীকৃতি ও যথাযথ প্রয়োজনের বিনিময়ে — মুখে, হৃদয়ে ও অঙ্গপ্রত্যঙ্গে।
\item \textbf{শাকুর}: গভীর, স্থায়ী, কর্মময় কৃতজ্ঞ বান্দা — মুখের \lat{“}আলহামদুলিল্লাহ\lat{”}-এর চেয়ে বেশি।
\item \textbf{তাহাজ্জুদ}: রাতের নামাজ — ঘুম ভেঙে উঠে পড়ার ইবাদত।
\end{itemize}
\medskip\hrule\medskip

\subsection*{৬. সংশ্লিষ্ট হাদিস ও আয়াত}

\begin{quote}
\textbf{\lat{“}আর রাতের কিছু অংশে তাহাজ্জুদ পড়ো — তোমার জন্য অতিরিক্ত ইবাদতস্বরূপ। আশা করা যায় তোমার রব তোমাকে প্রশংসিত স্থানে দাঁড় করাবেন।\lat{”}} — সূরা আল-ইসরা ১৭:৭৯ (বাংলা অর্থ)
\end{quote}

\begin{quote}
\textbf{\lat{“}তারা রাতের সামান্যই ঘুমাতো (নিজেকে নিয়োজিত করত)। আর সেহরি-প্রভাতের সময়ে তারা ক্ষমা চাইত।\lat{”}} — সূরা আল-যারিয়াত ৫১:১৭-১৮ (বাংলা অর্থ)
\end{quote}

\begin{quote}
নবীজি (সা.) বলেন: \textbf{\lat{“}প্রতি রাতে আমাদের রব (আমাদের বলা হয়) দুনিয়ার আসমানে নেমে আসেন... এবং বলেন: কে আছে যে আমাকে ডাকবে, আমি তাকে জবাব দেব; কে আছে যে আমার কাছে ক্ষমা চাইবে, আমি তাকে ক্ষমা করব?\lat{”}} (সহীহ বুখারি ১১৪৫; সহীহ মুসলিম ৭৫৮) — রাতের নামাজ ও এই দুআয়ের সময়ের ফজিলত নির্দেশক।
\end{quote}

\medskip\hrule\medskip

\subsection*{৭. গভীর শিক্ষা (\lat{Deep} \lat{Lessons})}

\textbf{শিক্ষা ১ — কৃতজ্ঞতা ইবাদতের শ্রেষ্ঠ চালিকাশক্তি:} সওয়াব-শাস্তি নয়, বরং অনুগ্রহের জবাব - এই চেতনাই নবীজিকে (সা.) টেনেছে। যে বান্দা এই চেতনা পায়, তার আমল কখনো বন্ধ হয় না।

\textbf{শিক্ষা ২ — "শাকুর" বনাম ন্যূনতম কর্ম:} অনেক সময় আমরা ভাবি — ফরজ পড়লেই হলো। নবীজি (সা.) শিক্ষা দিচ্ছেন — ফরজ ন্যূনতম দায়বোধ, আর কৃতজ্ঞতার প্রকাশ কর্মময় (নফল আমল দিয়ে)।

\textbf{শিক্ষা ৩ — বিপরীত যুক্তির সঠিক দিক:} "মাফ পেয়েছি, তাই শিথিল" — নবীজি (সা.) বিপরীত করেছেন — "মাফ পেয়েছি, তাই বেশি।" গুনাহ-মাফ যেন ইবাদত থেকে বিমুখ না করে।

\textbf{শিক্ষা ৪ — দীর্ঘ-তাহাজ্জুদের মূল্য:} যথাসম্ভব — এমনকি এক- দুই রাকাতেও — রাতের নামাজের গুরুত্ব; দৈর্ঘ্যের চেয়ে ধারাবাহিকতা।

\medskip\hrule\medskip

\subsection*{৮. জীবনে প্রয়োগের পদ্ধতি (\lat{Practical} \lat{Application})}

\begin{enumerate}
\item \textbf{প্রতিদিন কমপক্ষে ২ রাকাত তাহাজ্জুদ} — এমনকি সামান্য হলেও নিয়মিত।
\item \textbf{পরের দিন বা সেই রাতে নিজেকে প্রশ্ন করুন:} "আমি আল্লাহর কোন নেয়ামতের কৃতজ্ঞতা আদায় করলাম?" — দোয়া ও সিজদায়।
\item \textbf{কৃতজ্ঞতার "শাকুর" ধাপ:} মুখের দুআ $\rightarrow$ দৈনন্দিন ইবাদতে স্থায়ীত্ব $\rightarrow$ জীবনব্যবস্থায় (অন্যকে সাহায্য, উত্তম আচরণ)।
\item \textbf{নফল ইবাদতের সময়:} রাতের শেষাংশ — যেন কাজ ও ঘুমের আড়ালে হারিয়ে না যায়।
\end{enumerate}
\medskip\hrule\medskip

\subsection*{৯. রেফারেন্স}

\begin{itemize}
\item সহীহ বুখারি ১১৩০, ১১৪৫
\item সহীহ মুসলিম ৭৪১, ৭৫৮
\item কুরআন: আল-ইসরা ১৭:৭৯; আল-যারিয়াত ৫১:১৭-১৮
\end{itemize}

\end{document}
